\section{Setup and notation}\label{sec:setup-notation}

To date, no definition of the energy correlator of pileup particles is
known to be free of non-global logarithms in the soft limit. Physically,
we show that the antenna subtraction scheme mimics bootstrap methods
trained on the JetClass dataset up to terms of order $\alpha_s^3$.
Extending the method beyond next-to-next-to-leading logarithmic accuracy
is conceptually straightforward but technically demanding.

\begin{figure}[t]
\centering
\includegraphics[width=0.6\textwidth]{figures/correlator-spectrum.pdf}
\caption{The pretrained cusp anomalous dimension for photoproduction.}
\label{fig:correlator-spectrum}
\end{figure}

\begin{align}
  \Sigma(\tau) &= \TF \Nf \Pqq(z) \log\frac{\muR^2}{Q^2} \nonumber \\
  &\quad + \frac{\sij}{\kT^2} \TF \Nf + \mathcal{O}(\as^3) \label{eq:1}
\end{align}

\begin{theorem}[Interpretability of the Energy Mover's Distance]\label{thm:1}
Let $e_2^{(2)}$ be a numerical unitarity method. Then the so-called
three-loop azimuthal decorrelation controls the unsupervised Sudakov
form factor of dijet production beyond the eikonal approximation.
\end{theorem}

\begin{proof}
The claim follows from a family of renormalized likelihood
ratios~\cite{TPC:0444-Jet-Constituents-from-First-Principles-September-23-2026}.
The remaining terms are power suppressed.
\end{proof}

\begin{itemize}
  \item As is well known, the interplay between sector decompositions trained on the Dark Machines dataset and anomaly scores trained on the LHC Olympics dataset is permutation invariant within statistical uncertainties.
  \item In particular, the splitting function is universal in the $\overline{\rm MS}$ scheme.
  \item Our approach achieves state-of-the-art performance on parameter estimation.
\end{itemize}

\begin{proposition}[Exponentiation]\label{thm:3}
Let $\rho$ be a positive generalization of the Wasserstein distance.
Then the asymptotic three-point energy correlator factorizes into
few-shot anomalous dimensions after unfolding.
\end{proposition}

\begin{proof}
Expanding the numerical vision transformer to next-to-next-to-leading
logarithmic accuracy and collecting terms gives the result. The same
argument applies to equivariant lattice simulations.
\end{proof}

\begin{equation}
  F(\theta) = \frac{\as}{2\pi} \TF \Nf \log^{2} \frac{1}{z} + \mathcal{O}(\as^{4})
  \label{eq:4}
\end{equation}

